% !Mode: : "TeX: UTF-8"
%%% 本文档可自动生成研究生毕业论文的封面、摘要、目录、正文、参考文献、附录和致谢!等！
\documentclass[acaddeg]{Style/SCNUThesisV3} 
%% [2025/3/11   Version 3.0 South China Normal University, Thesis Template, Zhong Liuqiang]
%%% 参数可在acaddeg (学术学位academic degree) 和 profdeg (专业学位 professional degree) 中二选一
%%% 专业学位格式自动满足  华南师范大学教师教育学部 2024 年 11月颁布的华南师范大学研究生学位论文撰写规范！！
%%% 学术学位格式自动满足  华南师范大学数学科学学院 2024 年 11月颁布年颁布的 研究生学位论文撰写规范（适用学科：0701数学）！！
%%% 博士还是硕士, 需要将在 1Cover.tex 中输入.

\usepackage{Style/scnumacro}
%%% 自定义的宏包 scnumacro.sty, 已基本包含了所有的数学/表格等环境所需要的宏包



%%%%%%%%%%%%%%%%%%%%%%%%%%%%%%%%%%%%%图表字体和表头设置%%%%%%%%%%%%%%%%%%%%%%%%%%%%%%%%
\usepackage[labelsep=space]{caption}
%%% 自定义的宏包 scnumacro.sty, 已基本包含了所有的数学/表格等环境所需要的宏包

\renewcommand{\baselinestretch}{1.5} % 1.5 倍行距排版
%%图标标题字体设置
\DeclareCaptionFormat{myformat}{\textbf{#1 #2#3}} % #1 是标签，#2 是标题
\captionsetup[figure]{\wuhao format=myformat} % 应用到figure环境
\captionsetup[table]{\wuhao format=myformat} 
\renewcommand{\baselinestretch}{1.5} % 1.5 倍行距排版
%%%%%%%%%%%%%%%%%%%%%%%%%%%%%%%%%%%%%%%%%%%%%%%%%%%%%%%%%%%%%%%%%%%%%%%%%%%


%\hypersetup{colorlinks=false}%把链接变成红色
\hypersetup{colorlinks=true,
linkcolor=black,
filecolor=black,
urlcolor=black,
citecolor=black,}

\begin{document}
\let\cleardoublepage \clearpage
% 默认每一章应另起一页，每章首页应在奇数页面（右页）, 如果需要去掉空白页, 需要解开上述命令

%%% ===============================  
%%% ======  生成中英文封面  ====== 
% !Mode:: "TeX:UTF-8"
%%%论文封面
%%  教育硕士专业学位是按   华南师范大学教师教育学部 2024 年 6 月 6 日颁布的华南师范大学研究生学位论文撰写规范  所制定
%%  学术硕士专业学位是按 华南师范大学数学科学学院 2024 年 5 月 23 日颁布年颁布的 研究生学位论文撰写规范（适用学科：0701数学） 所制定

%%%%%%%%%%%%%%%%%%%%%%%%%%%%%%
 % 专业学位     中文封面内容  
%% “论文分类号”按《中国图书资料分类法》的分类号填写
%%% 具体可查询 http://ztflh.xhma.com
% 分类号 学术型:
% 计算数学 O24   应用数学 O29  
% 概率论与数理统计 O21   运筹学  O22
% 统计学 C8   课程与教学论  G427

\classification{}  % 分类号 
\UDC{} % UDC: 表示国际十进分类法(Universal Decimal Classification),是世界一般报告或者论文的封面在左上角注明分类号，便于信息交换和处理。
%%%  专业学位的 分类号、UDC：不填写
%%%  学术学位的  分类号: 有  计算数学 O24   应用数学 O29   概率论与数理统计 O21    运筹学  O22 统计学 C8   课程与教学论  G427
\confidential{公开}
%% “密级：学位论文密级分为公开、内部、涉密。内部论文和涉密论文严格按照学校审批的保密级别和保密期限填写，如内部*年、秘密*年；其他论文均填公开
\studentID{2020XXXXXX} 
%%   填写学位申请人的学号

 
\type{硕\ 士}      % 可选 硕士 或  博士
%\type{博\ 士}   
 
%%% ==================== 
%%% ====== 中文标题 ======
%%% ==================== 
%%% 论文标题: 应准确概括整篇论文的核心内容，简明扼要，一般不超过25个汉字。必要时可加副标题。
\title{\CASTunderline[\textwidth]{基于 Newman 错误分析的高中生数学解题}
\\ \vspace{0.18cm}
\CASTunderline[\textwidth]{错误研究 ---以三角恒等变换为例}} 
%% 如上是封面中的 论文中文标题, 需要人工分行

\titleab{基于 Newman 错误分析的高中生数学解题错误研究} 
\subtitle{}    % 中文副标题,  排版时自动居右, 这是可选项, 若不需要可注解掉


%%% ==================== 
%%% ====== 英文标题 ======
%%% ==================== 
%%% 英文标题 

\etitle{\CASTunderline[\textwidth]{Title: A Study of Mathematics Problem Solving }
\\ \vspace{0.18cm}
\CASTunderline[\textwidth]{Errors of High School Students Based on Newman}
\\ \vspace{0.18cm}
\CASTunderline[\textwidth]{Error Analysis --- Take the trigonometric constant}
\\ \vspace{0.18cm}
\CASTunderline[\textwidth]{transformations as an example}
} 

 
%%%==================== ==========  
 %%%==========  学位申请人 ==========  
\author{李建国} 
\eauthor{LI Jianguo}
%%% 英（外）文封面作者姓名按姓在前、名在后的格式书写，姓名需写全拼，如LI Jianguo


%%%==================== ==========  
%%%==========  学科专业  ==========  
%%% 学科专业: 申请人所在学科专业，必须严格按参照学科代码/专业学位代码字典中的二级学科名称/专业学位领域名称填写，若所选一级学科/专业学位类别下无二级学科/专业学位领域，则填写一级学科/专业学位类别名称，不可用简称。

%%%    学术学位   
%  \major{计算数学}
% \emajor{Computional Mathematics}
%  \major{应用数学}
% \emajor{Applied Mathematics}
%  \major{基础数学}   
% \emajor{Pure Mathematics}
%  \major{概率论与数理统计}   
% \emajor{Probability Theory and Mathematic Statistics}
%  \major{运筹学与控制论}   
% \emajor{Operations Research and Cybernetics}

%%%       专业学位
\major{教育硕士}   
\emajor{Master of Education}   %Major 仅出现在专业学位中 学科专业 的英文翻译, 时 需要,  学术学位不出现
%  \major{教育博士} % 专业学位名称
%\emajor{Doctor of education of Education}


%%%==================== ==========  
%%%==========  论文研究方向  ==========  
%%% 论文研究方向不能与学科专业名称完全相同，一个研究方向最多12个字，最多两个，以中文分号分隔，总数不超过24个字。
%%%%    学术学位: 研究方向 
%\specialization{微分方程数值解及应用} % 研究方向 
%%%%    专业学位: 
\specialization{学科教学(数学)}  
\especialization{Subject Teaching (Mathematics)} 
 

%%%==================== ==========
 %%%==========   所在学院  ========== 
 %%% 申请人所在二级培养单位的名称，必须填写与学校官网一致的规范名称，不可用简称。
\college{数学科学学院}  
\ecollege{School of Mathematical Sciences}   


%%%==================== ==========
 %%%==========  导师姓名及职称 ========== 
 %%%  填写指导教师的姓名与职称。若有两名指导教师，第一（校内）导师姓名和职称写在前面，第二（或校外）导师姓名和职称写在后面，中间空一格。第一（校内）导师必须与研究生管理系统一致。
\advisor{钟柳强} % 导师姓名
\ProfessionalTitle{教授} % 及职称
\eadvisor{ZHONG Liuqiang}  % 
%英（外）文封面作者姓名按姓在前、名在后的格式书写，姓名需写全拼，如LI Jianguo
\secondadvisor{ } %校外(行业)导师姓名, 仅在专业学位时需要,  学术学位不出现


%%%==================== ==========
 %%%==========  论文完成时间 ==========
 % 填写学位论文最终定稿打印成文的年月，封面的日期用汉字书写，如二〇二四年六月、二〇二四年十二月
\submityear{二〇二五} % 年, 用于中文封面
\submitmonth{五}         % 月, 用于中文封面
\submitengyear{2025}  % 年, 用于英文封面
\submitengmonth{MAY} % 月, 用于英文封面


\maketitle   % 生成中文封面
\makeengtitle   % 生成英文封面

%%% ===============================  
%%% ======  答辩合格证明  ======  提交电子版本需要纸质签名并扫描
\passthedefense

%%% ===============================  
%%% ======  中英文摘要及关键词  ====== 
% !Mode: : "TeX: UTF-8"
%%%(二)中英文摘要及关键词
%%%1．摘要是论文的内容不加注释和评论的简短陈述, 应做到不阅读全文就能获得与论文同等量的必要信息, 是一篇完整的短文, 可以独立使用、引用. 摘要应说明论文研究的目的、使用方法、结果和结论等. 
%%%2．中文摘要一般为1000字左右. 外文摘要为中文摘要的翻译, 所表述的内容与中文摘要一致. 外文摘要另起一页打印. 摘要中不用图、表、结构式公式及非公知公用的符号与术语. 
%%%3．关键词是为了文献标引工作而从论文中选取出来, 用以表示全论文主题内容信息的单词或术语, 一般为3—8个. 关键词字体加粗另起一行, 排在摘要的左下方, 按涉及的内容、从大到小排列. 多个关键词之间用分号隔开, 最后一个关键词不加标点, 应尽量用《汉语主题词表》等词表提供的规范词. 英文关键词、摘要内容应与本论文的中文关键词、摘要内容相对应. 
 

 %%%%%%% 中文摘要页: 

 %% 中文摘要内容
 \begin{cabstract}
《中国高考评价体系》是深化新时代高考内容改革的理论支撑和实践指南, 它指出学科素养是高考考查的关键内容之一. 与此同时, 逻辑推理素养作为高中数学六大核心素养之一, 不仅是数学严谨性的基本保证, 也是人们在数学活动中进行交流的基本思维品质. 因此, 如何基于高考评价体系进行逻辑推理素养的培养, 是目前一线教师重点关注的问题之一. 本文主要基于高考评价体系, 同时考虑《普通高中数学课程标准(2017年版2020年修订)》的要求, 以``圆锥曲线的方程"为切入点, 探索了高中生逻辑推理素养的培养策略, 在广东省汕尾市某高中进行了教学实验, 并检验了教学效果. 
\\
\indent
首先, 通过梳理相关文献, 界定了``推理与逻辑推理"和``逻辑推理素养"这两个核心概念. 然后, 分析了高考评价体系在数学学科的研究现状、逻辑推理素养的研究现状以及逻辑推理素养下圆锥曲线的教学现状. 发现虽然有较多学者关注逻辑推理素养的培养, 但关于考试与教学融合的研究工作尚不多见, 并且圆锥曲线与逻辑推理素养结合的研究较少.
\\
\indent
% 注: 在摘要环境中, 必须用上述两行来产生缩进的新段落
其次, 基于前面的文献研究, 结合高考评价体系, 对逻辑推理素养的水平划分框架与``四翼"的内涵进行了解读, 将逻辑推理素养与``四翼"均划分为三个水平等级, 水平越高意味着要求越高. 在此基础上, 选取四套新高考数学试卷, 以试卷中考查逻辑推理素养的试题为对象, 分析了其逻辑推理素养水平与``四翼"水平. 分析结果如下: (1) 试题对逻辑推理素养水平的考查主要集中在水平一与水平二, 且各卷在逻辑推理素养水平的命题情况基本一致;  (2) 试题对``四翼"的考查水平主要集中在基础性水平二与综合性水平二, 对应用性和创新性的考查较少. 
\\
\indent
然后, 依据《普通高中数学课程标准(2017年版2020年修订)》的要求编制了逻辑推理素养水平测试卷, 对广东省汕尾市某高中高二年级的学生进行了逻辑推理素养水平测试, 并以逻辑推理素养为主题, 对该校四位高二年级的数学教师进行了访谈. 利用SPSS分析了测试结果, 并总结了访谈结果, 得到以下结论: (1)大部分学生的逻辑推理素养水平处于水平一与水平二,极少学生能够达到水平三; 学生在合情推理方面的表现优于演绎推理; 学生的逻辑推理素养水平与学生的数学成绩呈正相关, 且男生与女生的逻辑推理素养水平不存在明显差异. (2)教师虽然无法准确地表述逻辑推理素养的内涵, 但能理解逻辑推理在数学学习中的重要性, 对圆锥曲线专题下培养学生逻辑推理素养的策略有一定思考; 虽然教师未对高考评价体系进行细致研究, 但对其核心内容都有一定的认识. 
\\
\indent
接着, 基于高考评价体系下的试卷分析结果, 结合高中生逻辑推理素养的水平调查与教师访谈, 提出了高考评价体系下高中生逻辑推理素养的培养策略: (1)注重直观与逻辑推理的辩证统一; (2)合理设置情境, 应对不同推理类型; (3)发展数学思维, 严谨逻辑表达; (4) 重视概念教学, 完善知识结构.
\\
\indent
最后, 依据制定的培养策略选取了实验班和对照班进行教学实验, 教学前后均对学生进行逻辑推理素养水平测试, 实验结果如下: (1)实验班学生的后测分数与前测分数几乎持平,对照班学生的后测分数相较于前测分数有所下降; (2)实验班在后测中达到水平三的人数比例高于前测; (3)在演绎推理方面, 实验班学生后测得分高于前测. 结果表明, 此次教学实验对学生逻辑推理素养的培养有一定的积极作用, 希望本文的研究能够为教学实践提供有益的参考. 
\end{cabstract}


 %% 中文关键词, 各关键词之间用逗号隔开, 最后一个关键词不加标点. 
 %% 关键词一般在3—8个之间. 
\begin{ckeywords}
逻辑推理素养, 高考评价体系, 圆锥曲线, 培养策略
\end{ckeywords}
 
 %%%%%%% 英文摘要页: 
 %%%%%%% 英文摘要内容和关键词应与本论文的中文关键词和摘要内容相对应. 
 
  %% 英文摘要内容
\begin{eabstract}	
The \textsl{`` China's College Entrance Examination Evaluation System"} is the theoretical support and practical guide for deepening the reform of the content of the new era college entrance examination. It points out that subject literacy is one of the key contents of the college entrance examination, and logical reasoning literacy, as one of the six core literacies in high school mathematics, is not only the basic guarantee of mathematical rigor but also the basic thinking quality for people to communicate in mathematical activities. Therefore, how to cultivate logical reasoning literacy based on the college entrance examination evaluation system is one of the key issues that frontline teachers currently focus on. Based on the college entrance examination evaluation system and considering the requirements of the \textsl{``National Curriculum Standard for General High School Mathematics (2017 Edition Revised in 2020)"}, this paper explores the cultivation strategy of logical reasoning literacy of high school students with the ``equation of conic curve" as the starting point, and conducts a teaching experiment in a high school in Shanwei City, Guangdong Province, and tests the teaching effect. 
\\
\indent
Firstly, by sorting out related literature, the core concepts of ``reasoning and logical reasoning" and ``logical reasoning literacy" were defined. Then, the research status of the college entrance examination evaluation system in mathematics, the research status of  logical reasoning literacy, and the teaching status of conic curve under logical reasoning literacy were analyzed. Although many scholars pay attention to the cultivation of logical reasoning literacy, there are few studies on the integration of testing and teaching, and there are few studies on the combination of conic curve and logical reasoning literacy. 
\\
\indent
Secondly, based on the previous literature research and the college entrance examination evaluation system, the level division framework and the ``four wings" connotation of logical reasoning literacy were interpreted, and both logical reasoning literacy and ``four wings" were divided into three levels, and the higher the level, the higher the requirements. On this basis, four sets of new college entrance examination mathematics papers were selected, and the logical reasoning literacy level and ``four wings" level of the questions on logical reasoning literacy in the papers were analyzed. The analysis results are as follows: (1) The examination of logical reasoning literacy level in the questions mainly focuses on level one and level two, and the proposition situation of the logical reasoning literacy level in each paper is basically the same; (2) The examination of ``four wings" level mainly focuses on foundational level two and comprehensive level two, and there is less examination of application and innovation. 
\\
\indent
Then, based on the requirements of the \textsl{``National Curriculum Standard for General High School Mathematics (2017 Edition Revised in 2020)"}, a test paper for logical reasoning literacy level was compiled, and a logical reasoning literacy level test was conducted on high school students in Grade 2 of a high school in Shanwei City, Guangdong Province, and four mathematics teachers in Grade 2 were interviewed on the topic of logical reasoning literacy. The test results were analyzed using SPSS, and the interview results were summarized. The following conclusions were drawn: (1) Most students' logical reasoning literacy level is at level one and level two, and very few students can reach level three; at the same time, their performance in conditional reasoning is better than deductive reasoning; there is a positive correlation between students' logical reasoning literacy level and their mathematical scores, and there is no significant difference in logical reasoning literacy level between male and female students. (2) Although teachers cannot accurately express the connotation of logical reasoning literacy, they can understand the importance of logical reasoning in mathematics learning and have some ideas on how to cultivate students' logical reasoning literacy under the topic of conic curve. Although teachers have not conducted detailed research on the college entrance examination evaluation system, they have a certain understanding of its core content. 
\\
\indent
Next, based on the analysis results of the college entrance examination evaluation system, the survey of high school students' logical reasoning literacy level, and the teacher interviews, the cultivation strategies for high school students' logical reasoning literacy under the college entrance examination evaluation system were proposed: (1) Emphasizing the dialectical unity of intuition and logical reasoning; (2) Reasonably setting up scenarios to deal with different types of reasoning; (3) Developing mathematical thinking and rigorous logical expression; (4) Attaching importance to concept teaching and improving knowledge structure. 
\\
\indent
Finally, based on the formulated cultivation strategies, an experimental class and a control class were selected for teaching experiments, and both classes were tested for their logical reasoning literacy level before and after the experiment. The results are as follows: (1) The post-test scores of the experimental class students are almost the same as the pre-test scores, while the post-test scores of the control class students have decreased compared to the pre-test scores; (2) The proportion of students who reach level three in the post-test of the experimental class is higher than that in the pre-test; (3) In deductive reasoning, the experimental class students score higher in the post-test than in the pre-test. The results show that this teaching experiment has a positive effect on cultivating students' logical reasoning literacy. It is hoped that this study can provide useful references for teaching practice. 
\end{eabstract}

 %% 英文关键词, 各关键词之间用逗号隔开. 外文关键词应与中文关键词相对应
\begin{ekeywords}
 logical reasoning literacy, college entrance examination evaluation system, conic curve, cultivation strategies.
\end{ekeywords} 
\makeabstract 

%%% ===============================  
%%% ======  目   录  ======
%%%      目录
\tableofcontents

%%%  以下内容选用 !!!
%%%% 插图索引
%\listoffigures
%%%% 表格索引
%\listoftables
\newpage\pagenumbering{arabic}\setcounter{page}{1}
%%% 论文正文开始页码编号重新由阿拉伯数字

%%% ===============================  
%%% ======  论文的页眉和页脚  ====== 
\pagestyle{fancy}
\fancyhf{}
\renewcommand{\headrulewidth}{0.8pt}  % 设置横线宽度 
%%% =====   奇数页为论文题目, 偶数页为华南师范大学硕士学位论文
\fancyhead[CO]{\titleinheader}  % 此处不包含中文封面的副标题, 这里可手工增加中文标题
\fancyhead[CE]{华南师范大学硕士学位论文}  % 可选  华南师范大学博士士学位论文
\cfoot{\thepage}
\fancypagestyle{plain}{
	\pagestyle{fancy}
}
%%% =========================== 
%%% ======　论文正文　===========
%%% 在撰写本部分之前, 一定要详细阅读华南师范大学学位论文基本要求
%%%%%% 正文已经默认小四号汉字, 可用如下命令做更改字号 \xiaosihao 

\mainmatter
%%%===============================================
%%%     ================ 　第 一章 　================
%%%===============================================
\chapter{绪论}\label{Cha: 1} %Cha1

%%%  学术学位论文第一章通常叫 前言,  专业学位论文 (profdeg) , 第一章通常叫 绪论,   
%%% 本部分属论文的主体部分, 主要介绍论文研究背景、论文所要解决的问题或研究的内容, 论文研究在教育工作中的实践价值和理论意义, 本研究领域国内外研究状况(或文献综述) , 论文的基本思路和论文结构等.

注: 学术学位论文的参考文献引用格式通常为 \cite{xu}.  专业学位论文的参考文献引用格式通常为 \upcite{xu}.

本章为绪论部分, 主要介绍论文的研究背景、论文将要解决的问题与论文内容、论文实践价值和理论意义、论文的研究方法以及论文结构. 
 \section{研究背景}
 目前, 由核心素养驱动的课程改革已经成为全球潮流, 我国在修订课程标准时, 对核心素养给予全面关注\upcite{xu}. 《中国高考评价体系》(下文简称高考评价体系) 也指出学科素养是高考内容的关键, 并且提出高考的核心功能为``立德树人"、``服务选才"与``引导教学"\upcite{中国高考评价体系}. 逻辑推理素养作为数学学科六大核心素养之一, 承担着重要的育人价值. 然而, 教师在教学实践中对逻辑推理素养的培养存在着一定的困难. 
 
\subsection{Test 研究背景}

\textbf{1. 基于高考评价体系对高中生逻辑推理素养的要求}

新高考数学命题框架的建构基础包括高考评价体系、高校人才选拔要求和国家课程标准\upcite{任子朝2019基于高考评价体系的数学科考试内容改革实施路径}, 其中高考评价体系是深化新时代高考内容改革的基础工程、理论支撑和实践指南\upcite{中国高考评价体系}. 高考评价体系的核心内容是``一核"、``四层"和``四翼", 而学科素养是``四层"考查内容的关键, 这是因为学科素养是学科必备知识、关键能力与情感态度价值观的综合体现\upcite{张开2019高考评价体系的研制解读}. 因此, 学科素养是新时代高考对高中生提出的要求. 

\textbf{2. 逻辑推理素养的育人价值}

数学是研究客观世界中数量关系和空间形式的科学, 通过逻辑推理、符号演算和科学计算认识世界\upcite{普通高等学校本科}.《普通高中数学课程标准(2017年版2020年修订) 》(下文简称《新课标》) 提出逻辑推理素养为数学核心素养之一, 并强调逻辑推理是得到数学结论、构建数学体系的主要方式, 是数学严谨性的基本保证, 是人们在数学活动中进行交流的基本思维品质 \upcite{教育部2020普通高中数学课程标准}. 史宁中在文\cite{史宁中2011数学的基本思想}中提出数学内部的发展依赖于逻辑推理. 因此, 培养学生的逻辑推理素养对于引导学生把握数学本质、促进学生思维发展以及塑造学生优秀的数学品质至关重要. 



\textbf{3. 逻辑推理素养的教学现状}

虽然《新课标》已经明确指出高中数学课程主要目标之一是发展学生核心素养, 但是在实际教学过程中仍然存在一些问题, 如教学手段和教学方式缺乏创新性, 高中数学的课程教学以教师的讲解为主, 学生在整个教学环节当中处于被动的学习状态 \upcite{肖玉2019新高考背景下高中数学教学模式的改革探析}. 另外受中国传统教育观念的影响, 演绎推理等价于逻辑推理的观念在教师和学生的头脑中根深蒂固. 这种观念导致教师和学生对归纳推理、类比推理的逻辑性充满怀疑, 在教学和学习实践中运用归纳推理、类比推理也会存在心理障碍\upcite{林玉慈2019高中数学课程中的逻辑推理及教学策略研究}. 喻平在文\cite{喻平2018数学核心素养的培养}中指出, 现在的教学偏重演绎推理, 让学生相信真理并无条件地接受真理, 这种局面不利于学生逻辑推理的发展. 因此, 正确理解逻辑推理素养的内涵, 对培养学生的逻辑推理素养至关重要.

2014 年, 《国务院关于深化考试招生制度改革的实施意见》出台, 对加强高考内容改革顶层设计提出要求, 明确指出要依据高校人才选拔要求和国家课程标准, 科学设计命题内容\upcite{2014国务院关于深化考试招生制度改革的实施意见}. 这为科学构建中国高考评价体系提出了明确目标, 提供了基本遵循\upcite{中国高考评价体系}. 高考评价体系明确规定了高考的功能之一是``引导教学", 因此, 高考评价体系与《普通高中课程方案》\upcite{教育部2020普通高中课程标准}共同对高考起指导作用.  本文将基于高考评价体系与《新课标》, 探索高中生逻辑推理素养的培养策略.

 \section{研究内容}
 本文主要基于以下问题: (1) 基于高考评价体系, 高考对逻辑推理素养的考查情况如何? (2) 高中生逻辑推理素养水平现状如何? (3) 高考评价体系下逻辑推理素养的培养策略与教学实验的效果如何? 
 
 针对研究问题, 本文聚焦以下三方面的研究内容: 
 
 (1) 结合《新课标》与高考评价体系, 解读逻辑推理素养的水平划分框架与``四翼"的内涵, 选取新高考数学试卷, 分析高考试卷中逻辑推理素养的考查情况.
 
 (2) 制定高中生逻辑推理素养水平测试卷与教师访谈提纲, 对高中生进行逻辑推理素养水平测试,并对教师进行逻辑推理素养相关访谈. 了解学生逻辑推理素养的水平现状与教师对逻辑推理素养的认知情况.
 
 (3) 基于理论与现状分析, 制定高考评价体系下高中生逻辑推理素养的培养策略, 依据培养策略进行教学实验, 并检验教学效果. 
 
 本文重点是基于高考评价体系与《新课标》, 分析高考试题, 了解高中生逻辑推理素养水平现状, 并探索日常教学中的逻辑推理素养培养策略, 从而促进高中生逻辑推理素养的发展.
 
 
 \section{研究意义}
教师在教学过程中大多从``教"的角度出发, 关注《新课标》对逻辑推理素养的教学要求. 但高考评价体系作为新时代高考内容改革的理论基础与实践指南, 明确指出高考的功能之一是引导教学. 因此, 在高中生的学习生涯中, 教学与高考密不可分. 本文将从考教结合的角度出发, 基于高考评价体系, 对高中生的逻辑推理素养培养进行理论与实践研究. 
  
 \textbf{1. 理论意义}
 
逻辑推理素养是数学六大核心素养之一, 学生在数学或者其它学科的学习中, 都不可避免地需要运用到逻辑推理思维.《新课标》虽然指出逻辑推理的主要形式, 但一线教师在理论上缺乏对逻辑推理素养内涵的认识, 本文从逻辑推理素养内涵出发, 对其进行解读, 丰富高中逻辑推理素养的理论内涵. 此外, 高考评价体系对高考和高中教育提出了新的评价要求, 本文基于高考评价体系, 探索高中生逻辑推理素养的培养策略, 有助于丰富高考评价体系在教学上的理论研究. 
 
\textbf{ 2. 实践意义}
 
目前, 我国对于学生逻辑推理素养的培养仍然存在一些问题, 如教学方式落后、教考结合不紧密等. 高考评价体系对高考考查内容、考查要求和考查载体进行了明确的说明, 它既是高考的指挥棒, 也对教学有一定的指导意义. 基于高考评价体系, 以圆锥曲线为切入点, 研究学生的逻辑推理素养培养, 不仅能够引导教学, 更能促进教学与高考之间的融合. 具体包括: (1) 通过调查高中生逻辑推理素养水平现状, 促进教师了解学生实际情况, 进而教师可以依据教学内容中的逻辑推理相关知识, 选择合适的教学方式, 做到因材施教. (2) 探索基于高考评价体系逻辑推理素养的培养策略, 以圆锥曲线为切入点, 将其应用于课堂实践教学, 加深教师对圆锥曲线与逻辑推理素养的理解, 为一线教师培养学生的逻辑推理素养提供借鉴与参考. 

 \section{研究方法}
 本文采用的研究方法主要有文献研究法、测验法、访谈调查法以及实验研究法.
 
 \textbf{1. 文献研究法}
 
 阅读相关文献, 整理目前基于高考评价体系的数学学科研究内容, 分析目前关于逻辑推理素养培养与测评的文章, 同时梳理逻辑推理素养下圆锥曲线的教学情况.
 
\textbf{2. 测验法}

基于理论分析与教师指导, 设计逻辑推理素养水平测试卷, 对高二年级的学生进行逻辑推理素养水平测试, 了解学生的逻辑推理素养水平. 
 
\textbf{3.访谈调查法}

选择高中数学教师作为访谈对象, 关注教师对逻辑推理素养内涵的了解程度、对逻辑推理素养培养策略的认知情况、对高考评价体系的了解与应用情况. 

\textbf{4.实验研究法}

选取两个班级进行教学实验, 一个班级为实验班, 另一个班级为对照班. 实验班依据作者制定的培养策略进行教学, 对照班按照常规方式进行教学. 通过对比两个班级逻辑推理素养测试的前后成绩, 分析教学实验对学生的逻辑推理素养培养是否有效. 
 
 本文的具体结构如图 \ref{fig:1.1} 所示: 
\begin{figure}[h]
	\centering
	\includegraphics[width=10.425cm, height=11.78cm]{论文框架.PNG}
	\caption{论文框架} %图片名称
	\label{fig:1.1}
\end{figure}
 
 \section{论文结构}
本文共分为七部分, 具体分布如下: 

第一章为本文的绪论, 主要介绍研究背景、研究问题与内容、研究意义、研究方法和路线以及本文结构等内容. 

第二章为文献综述与理论基础, 首先对本文需要使用的重要概念进行说明界定, 接着梳理研究内容的国内外研究现状, 进而凝练本文需要用到的教育理论. 

第三章为高考评价体系下逻辑推理素养考查情况分析, 首先基于《新课标》与高考评价体系对逻辑推理素养划分框架与``四翼"进行解读, 进而研究近两年新高考数学试卷逻辑推理素养的考查情况. 

第四章为高中生逻辑推理素养水平现状调查, 设计高中生逻辑推理素养水平测试卷与教师访谈提纲, 对学生进行测试并对教师进行访谈.

第五章结合新高考数学试卷的逻辑推理素养分析、高中生逻辑推理素养现状调查以及教师访谈结果, 制定高考评价体系下高中生逻辑推理素养的培养策略. 

第六章根据前文提出的培养策略进行教学实验. 为了验证制定的培养策略对提升高中生逻辑推理素养水平是否有效, 在广东省汕尾市某中学选择实验班与对照班进行为期一个月的教学实验, 并检测教学实验效果. 

第七章是本文的总结与展望, 得出本文的最终结论, 同时反思在研究过程中的不足, 对未来研究方向进行展望. 




%%%===============================================
%%%     ================ 　第  二 章 　================
%%%===============================================
% 当发现新的一章不是从奇数页开始时, 需要用如下命令插入新的空白页
\newpage
%\thispagestyle{empty}
% 这里可以放置额外的命令来确保完全是空白，例如：
\vspace*{\fill} % 这会尝试将空白填充到页面底部，但不是完美的解决方案。
\mbox{} % 使用 \mbox{} 创建一个完全不可见的盒子，确保页面是空白的。
\vspace*{\fill} % 再次确保页面底部也是空白的。
\newpage
%%%  如上命令插入新的空白页

\addtocontents{toc}{\vspace{-9pt}} % 在自动生成的目录中减少 9pt的垂直间距
\chapter{文献综述与理论基础}\label{Cha: 2}
\setcounter{equation}{0}
本章为文献综述与理论基础, 首先对本研究需要使用的重要概念进行界定说明, 接着梳理本研究内容的国内外研究现状, 进而凝练本研究需要使用的教育理论.

%正文是学位论文的核心部分, 一般由理论分析、研究设计(含研究方法、研究对象、理论假设、研究工具等) 与实施、数据资料、研究结果的统计分析、论证、结论等部分构成.其体例结构与写作形式可因学科专业特点和论文研究的内容、问题而不同.总体要求是: 观点明确, 实事求是、理论正确、逻辑清楚、层次分明、文字流畅、数据真实、公式计算推导结果无误.教育硕士学位论文一般为2—3万字.

 %\section{字体设置}

%中文字体很多, 但是常用的就那么几个.我们为 CJK 常用的六种中文字体定义了简单易用的命令

%宋体: $\backslash$songti, % 注, 源文件输入应该用  \songti

%黑体: $\backslash$heiti, % 注, 源文件输入应该用  \heiti

%仿宋: $\backslash$fangsong, % 注, 源文件输入应该用  \fangsong

%楷书: $\backslash$kaishu, % 注, 源文件输入应该用  \kaishu

%隶书: $\backslash$lishu, % 注, 源文件输入应该用  \lishu

%幼圆: $\backslash$youyuan % 注, 源文件输入应该用  \youyuan

%\cite{Alonso96: 385, NianFDonSg07Chn}
\section{核心概念界定}
本小节主要对本研究需要使用的重要概念进行界定说明, 具体有``推理与逻辑推理"与``逻辑推理素养"两个重要概念.

\subsection{推理与逻辑推理}

推理是思维的基本形式之一, 即推理是根据一个或一些判断而得出另一个判断的思维形式, 它是判断与判断的联结、过渡, 相当于语言中``因为"和``所以"之间的语句关系\upcite{刘江2004逻辑学}. 有 心理学专家在文\cite{朱智贤1989心理学大辞典}中提出: 推理过程是现代人逻辑思维的一种表现形式, 它是通过人类在头脑中对已有判断进行心理分析与判断综合, 从而引出新判断的一个心理过程. 不过近些年人们将推理的含义进行了延伸, 既包括推断、演绎、归纳和联系, 又包括猜想、实验和假设\upcite{鲍建生2009数学学习的心理基础与过程}. 在数学界中, 徐斌艳在文\cite{徐斌艳2001数学推理活动在数学教育中的意义}中将数学推理定义为: 人们在数学观念系统作用下, 由若干数学条件, 结合一定的数学知识、方法, 对数学对象形成某种判断的思维操作过程. 史宁中在文\cite{史宁中数学基本思想}中将数学推理定义为从一个数学命题判断到另一个数学命题判断的思维过程. 

 %%%===============================================
%%%     ================ 　第  三 章 　================
%%%=============================================== 
  
\addtocontents{toc}{\vspace{-9pt}}% 在自动生成的目录中减少 9pt的垂直间距
\chapter{高考评价体系下逻辑推理素养分析}\label{Cha: 3}
\setcounter{equation}{0}
高考评价体系是深化新时代高考内容改革的基础工程、理论支撑和实践指南, 高考试卷则是高考评价体系的实践, 基于高考评价体系的要求分析高考试题, 不仅有助于加深对高考评价体系的理解, 也有助于实现高考评价体系对教学的指导意义.本章选取四套新高考数学试卷, 分析试卷中逻辑推理素养的考查情况.


\section{新高考数学试卷逻辑推理素养水平分析}
本小节主要介绍逻辑推理素养水平的划分框架, 通过具体试题给出水平划分的实例, 进而分析四套新高考数学试卷逻辑推理素养水平的考查情况.

\subsection{逻辑推理素养水平划分}

作者以《新课标》中的逻辑推理素养水平划分为基础, 整理得出了逻辑推理素养水平划分标准, 具体标准见表  \protect \ref{table:3.1}.

\begin{table*}[h]
	\centering
	\caption{逻辑推理素养水平划分}
	\vspace{2mm}
	\begin{tabular}{cp{12cm}}
		\toprule
		\rule{0pt}{20pt}
		水平层次 &具体考查要求  \\
		\midrule
		\rule{0pt}{20pt}
		水平一& 1. 能够在熟悉的情境中, 用归纳或类比的方法, 发现数量或
		图形关系. 2. 能够在熟悉的数学内容中, 识别并通过例子理解不同推理的形式, 知道通过归纳推理、类比推理得到的结论是或然成立的, 通过演绎推理得到的结论是必然成立的. 了解熟悉的数学命题的条件与结论之间的逻辑关系; 掌握一些基本命题与定理的证明, 并有条理地表述论证过程. 3. 能够了解熟悉的概念、定理之间的逻辑关系. 4. 能够在交流过程中, 明确所讨论问题的内涵, 有条理地表达观点. \\
		\rule{0pt}{20pt}
		水平二 & 1. 能够在关联的情境中, 发现并提出数学问题, 用数学语言
		予以表达; 能够理解归纳、类比是发现和提出数学命题的重要
		途径. 2. 能够对与学过的知识有关联的数学命题, 通过对其条件与结论的分析, 探索论证的思路, 选择合适的论证方法予以证明, 并能用准确的数学语言表述论证过程; 能够通过举反例说明某些数学结论不成立. 3. 能够理解相关概念、命题以及定理之间的逻辑关系, 建立网状的知识结构. 4. 能够在交流的过程中, 围绕主题, 并进行有理有据的论述. \\
		\rule{0pt}{20pt}
		水平三&1. 能够在综合的情境中, 找到合适的研究对象并提出有意义的数学问题. 2. 能够掌握常用逻辑推理方法的规则, 理解其中所蕴含的思想. 对于全新的数学问题, 能够提出不同的假设前提, 推断结论, 形成数学命题. 对于较复杂的数学问题, 能够通过构建过渡性命题, 探索论证的途径, 解决问题, 并会用严谨的数学语言表达论证过程. 3. 能够理解建构数学体系的公理化思想. 4. 能够合理地运用数学语言和思维进行跨学科的表达与交流.\\
		\bottomrule
	\end{tabular} 
	\label{table:3.1}
\end{table*}

 

\textbf{例1}	 已知双曲线$C: \dfrac{{{x}^{2}}}{{{a}^{2}}}-\dfrac{{{y}^{2}}}{{{b}^{2}}}=1\left( a>0, b>0 \right) $的右焦点为$F (2, 0) $, 渐近线的方程为$y=\pm \sqrt{3}x$. 

(1) 求$C$的方程; 

(2) 过$F$的直线与$C$的两条渐近线分别交于$A$, $B$两点, 点$P ({{x}_{1}}, {{y}_{1}})$, $Q\left( {{x}_{2}}, {{y}_{2}} \right) $在$C$上,  过点$P$且斜率为$-\sqrt{3}$的直线与过点$Q$且斜率为$\sqrt{3}$的直线交于点$M$, 从下面\ding{172}\ding{173}\ding{174}中选取两个作为条件, 证明另外一个成立. 

\ding{172}$M$在$AB$上; \ding{173}$PQ//AB$; \ding{174}$\left| MA \right|=\left| MB \right|$.

\textbf{分析: }本题考查双曲线的几何性质、直线与双曲线的位置关系, 学生需要分析已知条件, 联立直线与双曲线, 选择合适的证明与计算思路, 推理得出点$M$的坐标, 从而再选择对应的条件, 推理证明所需要的答案. 本题考查了常见的双曲线、直线与点的关系, 需要学生对条件以及结论进行综合分析, 选择合适的思路进行推理证明, 因此属于逻辑推理水平二. 

\textbf{例2} \ 已知函数$f(x) $及其导函数${{f}^{'}}(x) $的定义域均为$R$, 记$g(x) ={{f}^{'}}(x) $.若$f(\dfrac{3}{2}-2x) $, $g(2+x) $均为偶函数, 则\\
A.$f(0) =0$\qquad  \qquad \quad B.$g(-\dfrac{1}{2}) $  \quad  \quad \qquad C.$f(-1) =f(4) $ \quad  \quad \qquad D.$g(-1) =f(2) $

\textbf{分析: }本题考查函数的性质, 解题的关键是通过$f(\dfrac{3}{2}-2x) $, $g(2+x) $均为偶函数的条件, 得出函数$f(x) $, $g(x) $的对称性以及周期性, 即满足《新课标》中对逻辑推理的定义: 由一些事实和命题出发, 依据规则推出其他命题的素养. 同时解题者需要建立函数奇偶性、对称性以及导数等概念之间的逻辑关系, 形成知识的网状结构, 因此该题属于逻辑推理素养水平二.

\begin{table*}  
	\begin{tabular}{c|p{4cm}p{4cm}p{4cm}}
		\toprule
		\rule{0pt}{18pt}
		类别  & 水平一 & 水平二& 水平三 \\
		\midrule
		\rule{0pt}{20pt}
		基础性&考查学科基本概念与原理& 调动单一的学科知识或技能解决问题& 运用单一学科基本思想方法分析问题、解决问题\\
		\rule{0pt}{20pt}
		综合性 & 考查多个核心知识点&考查不同层面的知识结构、能力网络结构 &在多项相互关联的活动组成的复杂情境下, 综合运用数学思想方法解决问题的能力\\
		\rule{0pt}{20pt}
		应用性 & 试题素材结合实际背景&  运用数学知识、思想和方法对实际	问题进行分	析研究、解决问题&设置挑战性的问题	情境、有跨度的问题引导学生自主探索, 掌握解决实际问题的方法\\
		\rule{0pt}{20pt}
		创新性 & 提供多种形式的材料、设计条件或结论开放解题	方法多样、答案	不唯一的试题& 创设新颖情	境, 创新试	题呈现方式&解答过程能够展现	考生分析问题、解决问题的思维过程\\
		\bottomrule
	\end{tabular} 
	\label{table:3.11}
\end{table*} \cite{xu}

由于高考评价体系对``四翼"中的应用性与创新性的要求较高, 因此, 作者在划分应用性以及创新性的水平时, 对应用性以及创新性的水平三的要求也较为严格. 为了更好地展现对``四翼"水平的划分情况, 下面以两个高考题目为例, 对``四翼"水平的具体划分过程进行说明.

  

%%%===============================================
%%%     ================ 　第 四章 　================
%%%===============================================%
\addtocontents{toc}{\vspace{-9pt}}% 在自动生成的目录中减少 9pt的垂直间距
\chapter{高中生逻辑推理素养水平现状调查}\label{Cha: 4}
\setcounter{equation}{0}
为了解高中生逻辑推理素养水平与教师对逻辑推理素养的认知情况, 作者依据《新课标》的要求编制了逻辑推理素养水平测试卷, 并设计了教师访谈提纲. 本章首先对学生进行测试, 利用$SPSS$ 软件分析测试结果, 了解学生的逻辑推理素养水平; 之后对教师进行访谈, 探讨高中生逻辑推理素养的培养现状.

\section{高中生逻辑推理素养水平测试与分析}
本节主要介绍逻辑推理素养水平测试的目的与对象、测试卷的具体编制过程以及测试结果分析.

\subsection{测试目的与对象}
对高中生进行逻辑推理素养水平测试的目的如下: 

(1) 了解高二年级学生的逻辑推理素养水平情况, 分析目前学生在逻辑推理方面存在的问题.

(2) 通过本次测试, 研究高二年级学生的逻辑推理素养水平是否与性别和数学成绩有关.

本次逻辑推理素养水平测试选择广东省汕尾市某中学高二年级的部分学生为对象, 该校高二年级共22个教学班级, 学生的成绩水平各不相同. 测试共选取了五个班级: 高二年级12班、13班、15班、21班与22班, 其中22班学生数学成绩较为优异, 12班与15班学生数学成绩处于年级中等水平, 21班学生数学成绩相对较差. 选择的测试班级中包含不同水平的班级. 因此, 作者认为选取这些班级进行测试有助于了解整个高二年级的逻辑推理素养水平.



 
%%%===============================================
%%%     ================ 　第 五 章 　================
%%%===============================================
\addtocontents{toc}{\vspace{-9pt}}% 在自动生成的目录中减少 9pt的垂直间距
\chapter{高考评价体系下高中生逻辑推理素养的培养策略}\label{Cha: 5}
\setcounter{equation}{0}
作者结合前面的文献梳理、高考评价体系下逻辑推理素养分析、学生逻辑推理素养水平现状调查以及教师访谈, 从教学实践的角度, 提出了高中生逻辑推理素养培养策略. 


\section{注重直观与逻辑推理的辩证统一}

从逻辑推理内部形式来看, 归纳推理与类比推理都需要学生经过观察、猜想得出答案. 在此过程中, 学生面对的材料可以是文字、图片、数字或者其它形式材料, 即学生需要通过感官直觉观察得出猜想, 再进行演绎推理验证. 因此, 在某些时候, 逻辑推理离不开直观感知.

 

%%%===============================================
%%%     ================ 　第 六章 　================
%%%===============================================
\addtocontents{toc}{\vspace{-9pt}}% 在自动生成的目录中减少 9pt的垂直间距
\chapter{高中生逻辑推理素养教学实验研究}\label{Cha: 6}
\setcounter{equation}{0}
为了验证培养策略的有效性, 作者在广东省汕尾市某中学高二年级选择了两个班级进行教学实验. 本章主要介绍实验的前期准备、实验的具体过程以及实验结果.
\section{教学实验准备}
本小节主要介绍教学实验目的与对象、教学实验中需要使用的材料等.

\subsection{实验目的与对象}
本次教学实验的主要目的是, 验证前文提出的培养策略对提升高中生逻辑推理素养水平是否有效. 教学实验的对象是广东省汕尾市某中学高二年级的21班与22班两个教学班, 其中22班为实验班, 使用作者依据培养策略制定的教学设计进行教学; 21班为对照班, 依旧由科任教师采取常规模式教学. 在实验前做出如下假设:高考评价体系下高中生逻辑推理素养的培养策略对学生逻辑推理素养水平有一定提升作用.

 


%论文中如有与导师或他人共同研究的成果, 必须明确指出; 如果引用他人的结论、观点、方法、数据或成果等必须明确注明出处, 并与参考文献一致, 严禁抄袭剽窃.论文一般采用篇末注(尾注) , 如有需要, 引文可用脚注.\footnote{示例: 这里可以加脚注}


 %%%===============================================
%%%     ================ 　第 七 章 　================
%%%===============================================
\addtocontents{toc}{\vspace{-9pt}}% 在自动生成的目录中减少 9pt的垂直间距
\chapter{总结与展望}\label{Cha: 7}
\setcounter{equation}{0}
本章为总结与展望, 首先总结了研究所得结论, 继而反思研究的不足之处, 最后给出了研究展望. 
\section{研究结论}
本文主要研究以下三个问题:  (1) 基于高考评价体系, 高考对逻辑推理素养的考查情况如何? (2) 高中生逻辑推理素养水平现状如何? (3)高考评价体系下逻辑推理素养的培养策略与教学实验的效果如何?

 针对问题(1), 作者基于高考评价体系, 以新高考数学试卷中考查逻辑推理素养的试题为对象, 分析了其逻辑推理素养水平与``四翼"情况. 得出以下结论: 对逻辑推理素养考查的水平主要集中在水平一与水平二, 对``四翼"的考查集中在基础性与综合性, 应用性与创新性方面涉及较少.
 
针对问题(2), 作者对广东省汕尾市某高中高二年级的学生进行逻辑推理素养水平测试, 并对该校教师进行访谈, 得出如下结论: 大部分学生处于逻辑推理素养水平一和水平二, 极少部分学生的逻辑推理素养水平达到水平三; 并且男生与女生之间的逻辑推理素养水平不存在显著性的差异; 学生逻辑推理素养水平与数学成绩呈正相关. 教师认为学生逻辑推理素养水平有待提高, 且在合情推理的表现优于演绎推理; 同时, 教师对逻辑推理素养的重要性有一定认识. 
 
%%%%%%%%%%%%%%%%%%%%%%%%%%%%%%%%%%%%%%
%%%%%%%%%%%%%%%%%%%%%%%%%%%%%%%%%%%%%%
\backmatter %结束章节自动编号
 \fancyhead[C]{  } 
 %%% 停止正文中页眉!
 \renewcommand{\headrulewidth}{0pt} 
%%%=============   参  考  文  献   =============%%%

\bibliographystyle{Style/SCNUbib}
\cleardoublepage
\phantomsection

\addtocontents{toc}{\vspace{-9pt}}% 在自动生成的目录中减少 9pt的垂直间距
\addcontentsline{toc}{chapter}{参考文献{\cdotfill}}

\bibliography{References/ref}
 
% 如下给出 Bib 数据库的路径和名称
%\bibliography{References/refs}
%\bibliography{/Users/zhongliuqiang/Math/Bibtex/0Student/Math}
%\bibliography{/Users/zhongliuqiang/Math/Bibtex/0Student/MathEdu}
%\bibliography{/Users/zhongliuqiang/Math/Bibtex/zhong}
%\bibliography{Math2021415}
%%% 注1: 为准确运行 bib 数据库, 对于本文档(Main.tex) , 需要 XeLaTeX 编译一次, BibTeX 编译一次, 再用 XeLaTeX 编译两次.
%%% 注2: 有必要时, 需要对得到的参考文进行手工微调.  
%%%  参考文献 应 小二号黑体居中，然后   隔行  编排正文主体内容
%%%注 3: 参考文一般格式 详见   0Readme.docx

\let\cleardoublepage\clearpage

\appendix
%%%=============　　　附　　　录　　　=============%%%
%%%(六) 附录
%%%  附录一般作为学位论文主体的补充项目, 主要包括, 供读者阅读的、不宜放在正中的、过分冗长的调查问卷、访谈提纲、量表、计算程序、统计推导等.附录要有编号编于正文后, 其页码与正文连续编排.
%%%% 注 1 同一个内容, 应该放在附录的同一章里, 但是可以在该章里下设小节!!!
%%% 不一定需要
\include{Data/3Appendix}



%%%=============　   　　致             谢　   　　=============%%%
%%%(八) 后记, 一般为致谢.
\include{Data/4Acknowledgement}


%%% =============================== 
%%%攻读硕士学位期间取得的研究成果及教学实践成果（必要时）
\achievement  % 不一定需要

%%% =============================== 
 %%% 学原创性声明和授权使用声明. 
\declare
 
\end{document}
